\subsection{Cognitive Foundations of Learning}
\subsubsection{Cognitive Load Theory (CLT)}

Human cognitive architecture is characterized by a working memory that has severe limitations in capacity and duration when processing novel information. Working memory typically holds only about seven chunks of information at a time and is restricted in its ability to simultaneously process and organize interacting elements. All information must first pass through this limited working memory before it can be stored in the unlimited long-term memory as permanent schemas. When instructional materials exceed these processing constraints, cognitive overload occurs, leading the learner to lose focus and stop processing the content deeply. Consequently, total cognitive load—the sum of intrinsic, extraneous, and germane loads—must be managed to ensure that the learner's cognitive capacity is not overwhelmed\cite{Cog2019Sweller}.

\paragraph{Intrinsic Load}
Intrinsic cognitive load refers to the inherent complexity of the material being learned, which is determined by its "element interactivity". This load is generated by the number of elements that a learner must attend to simultaneously to understand the subject matter. Because experts can use previously stored schemas in long-term memory to process multiple interacting squiggles or concepts as a single chunk, intrinsic load is always relative to the learner’s level of prior knowledge. It cannot be reduced by instructional design unless the learning objective itself is simplified or broken down. Therefore, a task that is high-load for a novice may impose minimal load on an expert due to their ability to activate existing knowledge structures in working memory\cite{perkins2024clt}.

\paragraph{Extraneous Load}
Extraneous cognitive load is unnecessary mental strain caused by the way information is presented or the activities a learner is required to perform. Poor design choices, such as forcing learners to mentally integrate spatially separated text and diagrams (the split-attention effect), increase this bad form of load. Extraneous load consumes working memory resources that could otherwise be used for meaningful learning, thereby hindering mastery and schema construction. Instructional interventions aim to minimize this load by clarifying wording, simplifying navigation, and avoiding redundant information. Effective system design ensures that the cognitive capacity is focused on the essential content rather than on managing a messy user interface or unclear instructions\cite{studiomg_cognitive_ai}.

\paragraph{Germane Load}
Germane cognitive load refers to the productive mental effort that learners expend to construct schemas and build a permanent store of knowledge. Modern conceptualizations view germane load as a redistributive function that moves working memory resources away from extraneous activities and toward the intrinsic aspects of the learning task. Well-structured instructional methods increase germane load by prompting learners to engage in reflection, retrieval practice, and self-explanation. This type of load is essential for deep understanding, as it represents the actual mental work required to automate skills and integrate new information with prior knowledge. By minimizing extraneous distractions, systems can maximize the cognitive resources available for this critical germane processing\cite{best2024Bjorn}.

\subsubsection{Handling Overly Complex Lessons}
To mitigate cognitive overload, the system employs dynamic chunking and segmenting to break high-intrinsic-load material into digestible, bite-size nodes. This "segmenting effect" ensures that learners have sufficient time and capacity to organize and integrate information from one portion of the lesson before the next segment begins. By breaking a continuous unit into learner-controlled segments, the system prevents the "zoning out" that happens when novel information is presented at a pace too fast for deep processing. This approach essentially creates interconnected nodes of knowledge, allowing the working memory to focus on building component models before attempting to understand the entire complex causal system\cite{mayer2003}.

The system also utilizes isolated elements sequencing, where individual interacting components are presented sequentially before being integrated into a fully complex form. This strategy is particularly effective for very high-element interactivity tasks that would otherwise exceed working memory limits if all interactions were taught at once. Once these isolated elements are stored in long-term memory, the learner only needs to learn the interactions between them, which significantly reduces the working memory load during the final integration phase. By managing the order of presentation, the system ensures that the intrinsic complexity of the subject matter is introduced at a rate that aligns with the recovery and limits of working memory resources\cite{Cog2019Sweller}.

Finally, aBridgeAI provides "fading support" or scaffolding to transition learners from studying worked examples to independent practice. Initially, the system provides fully solved examples to reduce extraneous load, as solving a conventional problem requires high-load strategies like means-ends analysis which leave little capacity for learning. As the learner's expertise increases, the system provides partially worked examples and completion tasks, eventually requiring independent problem solving. This adaptive approach avoids the "expertise reversal effect," where redundant guidance becomes an extraneous load for learners who have already acquired sufficient domain knowledge\cite{Renkl2003}.